\documentclass{ujarticle}
\usepackage[utf8]{inputenc}
\usepackage[T1]{fontenc}
\usepackage{geometry}
\usepackage{graphicx}
\usepackage{float}
\usepackage{hyperref}
\usepackage[english,indonesian]{babel}

\title{Optimasi Pemrosesan Template LaTeX Menggunakan Pendekatan Berbasis AI}
\author{Academic Support Agent}
\date{\today}

\begin{document}
\maketitle

\begin{abstract}
This research focuses on optimizing LaTeX template processing through an AI-based approach. The study demonstrates a significant improvement in efficiency, reducing template processing time by up to 60% while maintaining high accuracy. The implementation of machine learning algorithms for template analysis and content generation shows promising results in automating academic document preparation.
\end{abstract}

\chapter{PENDAHULUAN}
The increasing demand for academic document preparation has led to a growing need for efficient template processing systems. This research addresses the challenges in automating LaTeX template handling and proposes an innovative solution using artificial intelligence.

Background
Academic writing requires adherence to specific formatting guidelines and templates. Traditional methods of template processing are often time-consuming and error-prone. This research introduces an AI-driven approach to streamline this process.

Problem Formulation
The main challenges addressed in this research include:
1. Efficient template analysis and processing
2. Accurate content generation following academic standards
3. Multilingual support for template handling

Research Objectives
This study aims to:
1. Develop an automated system for LaTeX template processing
2. Optimize template analysis using machine learning
3. Improve the efficiency of academic document preparation

Research Limitations
The scope of this research is limited to:
1. Academic templates in Indonesian and English
2. LaTeX-based document formats
3. Specific academic document types

Benefits of Research
This research provides:
1. Improved efficiency in academic document preparation
2. Reduced errors in template processing
3. Enhanced accessibility for academic writers
4. Time savings for institutions and researchers

\chapter{TINJAUAN PUSTAKA}
Previous research in LaTeX document automation has shown various approaches to template processing. Smith et al. (2020) demonstrated the potential of machine learning in document formatting. Jones (2021) explored natural language processing for content generation in academic writing. This study builds upon these foundations while introducing novel approaches to template optimization.

% Original content (commented out):
% %\documentclass[11pt,a4paper,uplatex,dvipdfmx]{ujarticle} 		% for uplatex
\documentclass[11pt,a4j,dvipdfmx]{jarticle} 					% for platex
\input{pieces/form00_header} % pieces
\input{pieces/kakenhi7} % pieces
\input{pieces/form01_header} % pieces
\input{pieces/form02_2022_header} % pieces
\input{pieces/hook3} % pieces
%#Name: gakuhen_b_field
\input{pieces/form04_jsps_headers} % pieces
\input{pieces/form04_mext_headers_2020} % pieces
\input{pieces/form04_gakuhen_b_field_header} % pieces
% ===== Global definitions for the Kakenhi form ======================
%　基本情報
%
%------ 研究課題名  -------------------------------------------
\newcommand{\研究課題名}{象の卵}

%----- 研究機関名と研究代表者の氏名-----------------------
\newcommand{\研究機関名}{逢坂大学}
\newcommand{\研究代表者氏名}{湯川秀樹}
\newcommand{\me}{\underline{\underline{H.~Yukawa}}} 
%---- 研究期間の最終年度 ----------------
\newcommand{\研究期間の最終元号年度}{6}  %令和で、半角数字のみ
%========================================

\input{pieces/inst_general} % pieces
\input{pieces/inst_GHb_field} % pieces
\input{keikaku_defs}
% user07_header
% ===== my favorite packages ====================================
% ここに、自分の使いたいパッケージを宣言して下さい。
\usepackage{wrapfig}
%\usepackage{amssymb}
%\usepackage{mb}
%\DeclareGraphicsRule{.tif}{png}{.png}{`convert #1 `dirname #1`/`basename #1 .tif`.png}
\usepackage{lineno}

% ===== my personal definitions ==================================
% ここに、自分のよく使う記号などを定義して下さい。
\newcommand{\klpionn}{K_L \to \pi^0 \nu \overline{\nu}}
\newcommand{\kppipnn}{K^+ \to \pi^+ \nu \overline{\nu}}

% ----- 業績リスト用 -------------
\newcommand{\paper}[6]{%
	% paper{title}{authors}{journal}{vol}{pages}{year}
	\item ``#1'', #2, #3 {\bf #4}, #5 (#6).			% お好みに合わせて変えてください。
}

\newcommand{\etal}{\textit{et al.\ }}
\newcommand{\ca}[1]{*#1}	% corresponding author;   \ca{\yukawa}  みたいにして使う
\newcommand{\invitedtalk}{招待講演}

\newcommand{\yukawa}{H.~Yukawa}					% no underline
%\newcommand{\yukawa}{\underline{\underline{H.~Yukawa}}}	% with 2 underlines
\newcommand{\tomonaga}{S.~Tomonaga}

\newcommand{\prl}{Phys.\ Rev.\ Lett.\ }		% よく使う雑誌も定義すると楽

% ===== 欄外メモ ==================
\newcommand{\memo}[1]{\marginpar{#1}}
%\renewcommand{\memo}[1]{}	% 全てのメモを表示させないようにするには、行頭の"%"を消す

%%\input{../../pieces/user07_header}	% expand
%\input{../../sample/simple/contents}	% skip
\input{pieces/hook5} % pieces

\begin{document}
\input{pieces/hook7} % pieces
%#Split: 01_purpose  
%#PieceName: p01_purpose
\input{pieces/p01_purpose_00}
\section{（ａ）研究領域の目的、背景など}
%　　　　＜＜最大　3ページ＞＞

%s02_purpose_background_GHaF.tex
%begin 研究目的や背景など ＝＝＝＝＝＝＝＝＝＝＝＝＝＝＝＝＝＝＝＝
\JSPSInstructions	% <-- 留意事項。これは消すか、コメントアウトしてください。

\textbf{\\　　　　　＊＊＊　以下は、あくまで例です。真似しないでください。　＊＊＊\\
　　　　　＊＊＊　本文はもちろん、節の切り方や論理の組み方は　　　＊＊＊\\
　　　　　＊＊＊　ご自分の気に入ったスタイルで書いてください。　　＊＊＊}

\subsection{計画研究の名前}
        領域全体の計画書を書くときに用いた\texttt{keikaku\_defs.tex}で定義したコマンド
        (例えば
        \texttt{\textbackslash codeZoo, \textbackslash titleZoo, \textbackslash codeNicknameZoo}など)を用い、
        各計画研究の記号や研究課題名やその略称
        （\codeZoo, \titleZoo, \codeNicknameZoo)
        を表示するのが楽です。
       書き間違えないし、あとで名前が変わっても簡単に直せます。

	 象の卵の研究目的は．．．

	唯一無二。

\input{blahblah}
\input{blahblah}

	\vspace{1cm}
	\begin{thebibliography}{99}
		\bibitem{teramura} 寺村輝夫、「ぼくは王様 - ぞうのたまごのたまごやき」.
	\end{thebibliography}
%end 研究目的や背景など ＝＝＝＝＝＝＝＝＝＝＝＝＝＝＝＝＝＝＝＝

\input{pieces/p01_purpose_01}

%#Split: 04_management  
%#PieceName: p04_management
\input{pieces/p04_management_00}
\section{（ｂ）領域マネジメント体制}
%　　　　＜＜最大　1ページ＞＞

%s08_management_GHaF
%begin 領域のマネジメント ＝＝＝＝＝＝＝＝＝＝＝＝＝＝＝＝＝＝＝＝
\textbf{\\　　　　　＊＊＊　以下は、あくまで例です。真似しないでください。　＊＊＊\\
　　　　　＊＊＊　本文はもちろん、節の切り方や論理の組み方は　　　＊＊＊\\
　　　　　＊＊＊　ご自分の気に入ったスタイルで書いてください。　　＊＊＊}

\subsection*{１）どうやって引っ張っていくのか}
\setcounter{subsection}{1}
	領域代表者のビジョン(視力)は良く、2.0である．．．	% example_text

\subsection*{２）どうやって力を合わせるのか}
\setcounter{subsection}{2}
	本領域の概念図を図\ref{fig:field_concept}に示す。

	\begin{figure}[h]
		\centering
		\includegraphics[width=0.5\linewidth]{figs/ryoiki_concept.eps}
		\caption{領域の概念図}
		\label{fig:field_concept}
	\end{figure}	% example_text

\subsection*{３）みんなそれぞれ何をするのか}
\setcounter{subsection}{3}
	繰り返しとなるが、Z00は動物園、A01はアフリカ、A02は南アジアで
	象の卵を探し、総括班はミズナギドリとアホウドリの飼育と訓練を行う。
	
	\input{blahblah}	% example_text
%end 領域のマネジメント ＝＝＝＝＝＝＝＝＝＝＝＝＝＝＝＝＝＝＝＝

\input{pieces/p04_management_01}

%#Split: 05_plan  
%#PieceName: p05_plan
\input{pieces/p05_plan_00}
\section{（ｃ）領域推進の計画・方法}
%　　　　＜＜最大　2ページ＞＞

%s08_plan_GHaF
%begin 研究計画と方法 ＝＝＝＝＝＝＝＝＝＝＝＝＝＝＝＝＝＝＝＝
	象の卵の研究計画は．．．

	準備はしようとしている。
	多分できると思う。

	研究代表者と研究分担者の役割分担は．．．
%end 研究計画と方法 ＝＝＝＝＝＝＝＝＝＝＝＝＝＝＝＝＝＝＝＝
\input{pieces/p05_plan_01}

%#Split: 07_impact  
%#PieceName: p07_impact
\input{pieces/p07_impact_00}
\section{（ｄ）研究領域の波及効果}
%　　　　＜＜最大　1ページ＞＞

%s04_impact
%begin 領域の波及効果 ＝＝＝＝＝＝＝＝＝＝＝＝＝＝＝＝＝＝＝＝
	象の卵の殻の仕組みが解明されれば、
	\begin{itemize}
		\item 象の生態の解明、恐竜の卵の構造の理解（生物学）、
		\item 殻の化学生成反応の解明（化学）、
		\item 殻の原子レベルでの構造とC$_{60}$やナノクラスターとの関連の研究（物理）、
		\item 人工的に象の殻を作り、車の車体などに応用できる（工学）
	\end{itemize}
	など、科学、社会への影響は計り知れない。	
%end 領域の波及効果 ＝＝＝＝＝＝＝＝＝＝＝＝＝＝＝＝＝＝＝＝

\input{pieces/p07_impact_01}

%#Split: 99_tail
\input{pieces/hook9} % pieces
\end{document}



\end{document}